% Capítulo incluido desde KalaScript.tex.
\section{Base de investigación}

Reúno aquí las fuentes que tomo como referencia para el proyecto. Me interesan
sus aportes sobre escritura, notación, composición y legibilidad. Esta bibliografía
recoge secciones relevantes; no supone una revisión exhaustiva de cada obra ni
de todas las disciplinas que quiero abordar.

\subsection{Referentes propuestos}

\begingroup\small
\begin{longtable}{@{}>{\raggedright\arraybackslash}p{\dimexpr 0.18\linewidth-1.3333333333333333\tabcolsep\relax}>{\raggedright\arraybackslash}p{\dimexpr 0.32\linewidth-1.3333333333333333\tabcolsep\relax}>{\raggedright\arraybackslash}p{\dimexpr 0.5\linewidth-1.3333333333333333\tabcolsep\relax}@{}}
\toprule
\textbf{Referente} & \textbf{Fuente} & \textbf{Aplicación y límite} \\ \midrule
\endfirsthead
\toprule
\textbf{Referente} & \textbf{Fuente} & \textbf{Aplicación y límite} \\ \midrule
\endhead
\bottomrule\endfoot
Hangul & \href{https://www.hangeul.go.kr/en/exhi/dailyExhibition.do}{National Hangeul Museum} & Relaciones sistemáticas entre formas. Sus mecanismos fonéticos no se trasladan automáticamente a conceptos. \\[4pt]
Shavian & \href{https://www.unicode.org/versions/Unicode17.0.0/core-spec/chapter-8/}{Unicode 17, capítulo 8} & Contrastes de altura y pares relacionados. No garantizan una silueta única para cada palabra. \\[4pt]
Blissymbolics & \href{https://www.blissymbolics.org/index.php/component/content/article?Itemid=0\&id=10}{BCI: composición y caracteres clave} & Reutilización y vocabulario ampliable mediante convenciones. La composición requiere definiciones compartidas. \\[4pt]
Gregg & \href{https://greggshorthand.github.io/anunit01.html}{Manual, unidad 1} & Enlaces, curvas, longitud y formas abreviadas. La rapidez real depende del aprendizaje y de la legibilidad. \\[4pt]
Grafonomía & \href{https://graphonomics.net/}{International Graphonomics Society} & Estudio experimental del proceso de escritura y otras habilidades gráficas. \\[4pt]
Dutton & \href{https://kennaway.org.uk/conlang.dir/Speedwords.dict}{Diccionario transcrito por Richard Kennaway} & Afijos de oposición diferenciados de otros negativos; no implica que todos los conceptos formen pares. \\[4pt]
Peirce & \href{https://gnusystems.ca/Lowell2.htm}{Lowell Lecture 2, manuscrito de 1903 transcrito} & Negación y alcance espacial. KalaScript usa un marcador explícito de negación; no implementa los grafos de Peirce. \\[4pt]
RDF & \href{https://www.w3.org/TR/rdf12-concepts/}{RDF 1.2 Concepts}, \href{https://www.w3.org/TR/swbp-n-aryRelations/}{relaciones n-arias} & Entidades, relaciones y argumentos definidos. El esquema no garantiza por sí mismo comprensión universal. \\[4pt]
APL / Iverson & \href{https://www.jsoftware.com/papers/tot1.htm}{Notation as a Tool of Thought, 1980} & Regularidad de operaciones, abstracción y capacidad de sugerir expresiones relacionadas. La compacidad no es el único objetivo. \\[4pt]
\end{longtable}
\endgroup

\subsection{Fundamentos complementarios}

\begin{itemize}
\item \href{https://wals.info/chapter/s1}{WALS: organización de la gramática}, \href{https://wals.info/chapter/141}{sistemas de escritura}: separar escritura, estructuras lingüísticas y significado.
\item \href{https://wals.info/chapter/77}{WALS: evidencialidad}: distinguir origen de evidencia de confianza en una afirmación.
\item \href{https://aclanthology.org/J14-1005/}{Kuhn, 2014: A Survey and Classification of Controlled Natural Languages}: precisión, expresividad, naturalidad y simplicidad como dimensiones, no como puntuación de calidad.
\item \href{https://lean-lang.org/theorem_proving_in_lean4/Dependent-Type-Theory/}{Lean: teoría de tipos}: tipos, términos y funciones como referencia de precisión formal.
\item \href{https://www.w3.org/TR/mathml4/}{MathML} y \href{https://openmath.org/standard/om20-2019-07-01/omstd20.html}{OpenMath}: distinguir presentación y significado, con definiciones de símbolos extensibles.
\item \href{https://www.w3.org/TR/skos-primer/}{SKOS}: conceptos, etiquetas, definiciones y notas de alcance.
\item \href{https://www.w3.org/TR/prov-primer/}{PROV}: procedencia de información y procesos que producen resultados.
\item \href{https://www.bipm.org/en/publications/si-brochure}{BIPM: SI}, \href{https://jcgm.bipm.org/vim/en/2.26.html}{VIM: incertidumbre}: cantidades, unidades e incertidumbre de medida.
\item \href{https://sbml.org/documents/what-is-sbml/}{SBML}: representación de modelos biológicos mediante núcleo y extensiones especializadas.
\item \href{https://www.unicode.org/faq/private_use.html}{Unicode: caracteres de uso privado}: posibilidad de codificar una escritura propia, con documentación y fuentes acordadas.
\end{itemize}

\subsection{Evidencia sobre legibilidad y aprendizaje}

\textbf{Miton y Morin, 2021.} \href{https://pmc.ncbi.nlm.nih.gov/articles/PMC8346946/}{Graphic complexity in writing systems}. Análisis de más de 47 000 caracteres de 133 escrituras. La complejidad está relacionada con el tipo de escritura; el trabajo no establece un número universal de trazos óptimo para signos nuevos.

\textbf{Zhang, Wang y Dai, 2026.} \href{https://link.springer.com/article/10.1007/s11145-026-10837-x}{The effect of motor sequence on written word production of a new script moderated by script experience and visual complexity}. Estudio con 78 participantes y copia diferida de caracteres de Hangul desconocidos. La presentación animada de trazos tuvo efectos diferentes sobre inicio, precisión y ejecución. No es una prueba del aprendizaje de un lenguaje completo.

\subsection{Estado de las especificaciones consultadas}

MathML 4 se consultó como \textbf{Working Draft del 4 de junio de 2026}. RDF 1.2 Concepts se consultó como \textbf{Candidate Recommendation Snapshot del 7 de abril de 2026}. No se presentan como recomendaciones finales. La página del SI enlazaba la novena edición actualizada en 2026.

Las decisiones de KalaScript son propuestas de diseño derivadas de esta investigación; ninguna de estas fuentes valida por sí sola los signos del prototipo.

\subsection{Consulta complementaria para la ampliación}
Se consultaron fuentes primarias en internet para precisar las nuevas reglas.
Las formas de los signos y la selección de funciones son decisiones propias de
KalaScript, pendientes de evaluación de uso.
\begin{itemize}
\item \href{https://www.w3.org/TR/owl2-primer/}{W3C, OWL 2 Primer, segunda edición
(Recomendación de 2012)}, secciones sobre clases complejas y restricciones:
referencia para separar individuos, clases y composiciones con condiciones.
KalaScript no implementa OWL ni sus motores de inferencia.
\item \href{https://www.cs.cmu.edu/~fp/courses/15816-s10/lectures/19-PDL.pdf}{André
Platzer, notas sobre lógica dinámica proposicional, CMU, 2010}, definición de
estructura y accesibilidad: referencia para la interpretación de posibilidad.
Se adopta aquí un fragmento proposicional con modelo explícito.
\item \href{https://15316-cmu.github.io/2020/lectures/03-safety.pdf}{Matt Fredrikson,
Semantics, Safety, \& Dynamic Logic, CMU, 2020}, definición de programas y
ejecución: referencia para asignación, secuencia y comprobación del bucle antes
del cuerpo. El modelo entero de KalaScript y su vocabulario están delimitados
en la extensión de gramática.
\item \href{https://ftp.cs.ucla.edu/pub/stat_ser/r350.pdf}{Judea Pearl, Causal
inference in statistics: An overview, 2009}: referencia para distinguir
intervención de observación. R6 se restringe a una influencia en modelos
deterministas acíclicos; no es un procedimiento de descubrimiento causal.
\end{itemize}
